% limbo.tex — the Limbo field review (compiled manuscript).
%
% Venue: Physics and Imaging in Radiation Oncology (phiRO), Elsevier / ESTRO
%   (ISSN 2405-6316). Article type: Review article. Open access.
% Class : elsarticle.cls — Elsevier's official article class (LPPL 1.3+). It IS on CTAN
%   and tectonic can fetch it, but it is vendored here (elsarticle.cls + the numbered
%   Vancouver bibliography style elsarticle-num.bst) so the build is self-contained and
%   the exact class/style version is pinned; see ELSEVIER_CLASS_PROVENANCE.md.
%   Build: tectonic limbo.tex (gated on verify_citations.py; see build.sh / reproduce.sh).
%
% This review reports NO new measurement and NO number of its own. Every empirical
% statement is attributed to a verified external source in limbo.bib / CITATIONS.md;
% verify_citations.py fails the build on any phantom \cite or unresolved identifier.
\documentclass[review,number]{elsarticle}

\usepackage{lmodern}        % scalable Latin Modern shapes (silences CM 5.5pt size substitutions)
\usepackage{amsmath,amssymb}
\usepackage{booktabs}
\usepackage{microtype}
\usepackage{url}            % catcode-safe \url for the JMLR/PMLR proceedings URLs (underscores)
\usepackage{lineno}\linenumbers   % phiRO requires line numbers in the submitted manuscript
\emergencystretch=2em

% elsarticle loads natbib itself (numbered scheme via the [number] class option); add
% sort&compress to match an ascending, range-compressed numbered rendering.
\biboptions{sort&compress}

\newcommand{\Dstar}{D^{*}}
\newcommand{\Ktrans}{K^{\mathrm{trans}}}

\journal{Physics and Imaging in Radiation Oncology}

\begin{document}

\begin{frontmatter}

\title{Trustworthy uncertainty quantification for quantitative and diffusion body MRI,
and its decision-use in MR-guided adaptive radiotherapy: a field review}

% Solo author; corresponding email is institutional (ak5232@columbia.edu), matching the
% cover letter. ORCID carried as an author footnote (elsarticle has no \orcid macro).
\author[col,cub,asi]{Avery Karlin\corref{cor1}\fnref{orcid}}
\ead{ak5232@columbia.edu}
\cortext[cor1]{Corresponding author}
\fntext[orcid]{ORCID:\ 0000-0003-3848-6782}

\affiliation[col]{organization={Department of Applied Physics and Applied Mathematics, Columbia University}, city={New York}, state={NY}, country={USA}}
\affiliation[cub]{organization={Department of Computer Science, University of Colorado Boulder}, city={Boulder}, state={CO}, country={USA}}
\affiliation[asi]{organization={Annealing Signet Institute}, city={New York}, state={NY}, country={USA}}

\begin{abstract}
In MR-guided adaptive radiotherapy, a quantitative or diffusion MRI biomarker is
increasingly offered not as an image but as a \emph{number} that can steer dose --- a
dose-painting prescription, or a response-adaptive change to the plan. Acting on that
number means acting on the \emph{uncertainty} attached to it, and the question this
review surveys is whether that uncertainty is trustworthy enough to drive a dose change.
We read the field as the path a biomarker error bar must travel to reach the dose
decision: \emph{trust} (is the bar honest --- calibrated and covering?), \emph{value}
(does a trusted bar change the decision?), and \emph{action} (what happens when it is
deployed to drive dose on an MR-linac?), over a foundations layer (what is measured, and
with what estimator) and beneath a map of open problems. We find that trust is
conditional: broadly defensible for the diffusion coefficient $D$ and ADC, but fragile
for the perfusion parameters $f$ and $\Dstar$, whose measured test--retest reliability is
poor and for which deep-learning uncertainty quantification typically reports only
marginal, aggregate calibration. The decision-analytic apparatus to price a calibrated
interval --- net benefit, decision calibration, value of information --- is mature but
rarely connected to quantitative-MRI uncertainty; and the propagation of \emph{biomarker}
uncertainty into the adaptive dose decision remains the immature link. We close with four
open problems for biomarker-driven adaptive radiotherapy, sited where trust, value, and
action meet. This review reports no new measurement; its contribution is synthesis,
taxonomy, and gap identification.
\end{abstract}

\begin{keyword}
MR-guided adaptive radiotherapy \sep dose painting \sep quantitative MRI \sep
intravoxel incoherent motion \sep uncertainty quantification \sep value of information
\end{keyword}

\end{frontmatter}

\section{The dose decision the biomarker must serve}
\label{sec:scope}

Adaptive MR-guided radiotherapy increasingly proposes to steer dose from what a
quantitative or diffusion MRI sequence \emph{measures}, not only from where the anatomy
sits. On an MR-linac the plan can be re-optimised between or within fractions, and an
emerging paradigm would let a functional or quantitative biomarker --- an apparent
diffusion coefficient (ADC), an intravoxel-incoherent-motion (IVIM) parameter, a
perfusion measure --- raise or lower dose to a sub-volume according to estimated biology
\cite{otazo2021,hall2022}. The biological premise is \textbf{dose painting}: prescribing
a non-uniform dose from a target volume derived from functional or molecular images
\cite{ling2000,bentzen2005,bentzen2011}. What such a prescription actually consumes is
not the biomarker but a \emph{number with an error bar}; and the unsolved link is whether
that error bar is honest and stable enough to justify changing a dose.

This review surveys that link. We read the literature as the path a biomarker error bar
must travel before it can reach a dose decision:

\begin{enumerate}
  \item \textbf{Trust} --- is the error bar honest (calibrated, covering)?
  \item \textbf{Value of information} --- given a trusted bar, does it change the decision?
  \item \textbf{Action} --- when deployed to drive dose on an MR-linac, where does it hold and where does it break?
\end{enumerate}

\noindent over a \textbf{foundations} layer (what is measured, and with what estimator)
and beneath a \textbf{gap map} of the field's open problems. Because the dose decision is
the organizing concern, we foreground the action setting
(Sec.~\ref{sec:action}) first, then walk foundations
(Sec.~\ref{sec:foundations}), trust (Sec.~\ref{sec:trust}), and value
(Sec.~\ref{sec:voi}) as the question of whether the quantitative-MRI side can supply what
that decision demands. The survey covers uncertainty quantification (UQ) for
quantitative/diffusion body MRI and its decision-use in MR-guided adaptive RT; it
excludes non-quantitative MRI, neuro-only work, generic computer-vision UQ, and
non-MR-guided RT (the scope boundary is recorded in the repository's
\texttt{ASSUMPTIONS.md}). Table~\ref{tab:axis} lays out the survey axis and the thematic
buckets it induces.

\begin{table}
\caption{The survey axis. The field is classified by the question each work answers
about \emph{somebody's} uncertainty estimate, foregrounding the action setting --- the
dose decision --- that the other axes serve; a cross-cutting gap map
(Sec.~\ref{sec:gaps}) sits over the axes at the seams between them.}
\label{tab:axis}
\centering
\begin{tabular}{p{0.20\linewidth}p{0.68\linewidth}}
\toprule
axis & what the bucket asks \\
\midrule
Action & When deployed to drive dose on an MR-linac, where does a biomarker's error bar hold and break? (Sec.~\ref{sec:action}) \\
Foundations & What is the estimand, and what estimator (with what uncertainty) produced it? (Sec.~\ref{sec:foundations}) \\
Trust & Is the error bar honest --- calibrated and covering? (Sec.~\ref{sec:trust}) \\
Value of information & Given a trusted bar, does it change an oncologic dose decision? (Sec.~\ref{sec:voi}) \\
Gap map & Where, at the seams between the axes, do the field's UQ-trust questions remain open for adaptive RT? (Sec.~\ref{sec:gaps}) \\
\bottomrule
\end{tabular}
\end{table}

\section{The action setting: biomarker-driven adaptive radiotherapy}
\label{sec:action}

The action setting where a quantitative-MRI biomarker most directly drives dose is the
MR-Linac. Integrating diagnostic-quality MRI with a linear accelerator
\cite{lagendijk2008,raaymakers2009} matured into a clinically feasible 1.5\,T MRI-guided
system \cite{raaymakers2017,lagendijk2014} and an emerging paradigm for adaptive
radiation oncology \cite{otazo2021,hall2022}. The biological premise is \textbf{dose
painting} --- prescribing a non-uniform dose from a biological target volume derived from
functional/molecular images
\cite{ling2000,bentzen2005,bentzen2011}% QUOTE-3 RESOLVED (ling2000) — verbatim re-pull, Ling et al. 2000, IJROBP 47(3):551-560, Abstract/Results (PMID 10837935): "Incremental to the concept of gross, clinical, and planning target volumes (GTV, CTV, and PTV), we propose the concept of 'biological target volume' (BTV) and hypothesize that BTV can be derived from biological images and that their use may incrementally improve target delineation and dose delivery." Paraphrase above is faithful; no number used. See CITATIONS.md.
{} --- with DWI/ADC a leading candidate response biomarker \cite{leibfarth2018} and
quantitative MRI proposed to update the plan to observed tumour response
\cite{vanhoudt2021,otazo2021}. Clinical translation is real but uneven: a randomised
focal boost to an MRI-defined intraprostatic lesion improved biochemical control in the
FLAME trial \cite{kerkmeijer2021} and dose-painting-by-numbers has been delivered in
early-phase studies \cite{berwouts2013}, yet a phase~III \textsuperscript{18}F-FDG-PET
dose-redistribution trial found no overall tumour-control benefit \cite{deleeuw2024} ---
so steering dose from an image map is supported in principle but not settled in outcome.
Diffusion biomarkers are now being tested on the MR-linac itself, e.g.\ longitudinal
per-fraction DWI during pancreatic stereotactic RT as a prognostic marker
\cite{bisgaard2024}.

Two cautions make the dose decision conditional on the biomarker's uncertainty, not just
its value. First, \textbf{uncertainty must be carried into the optimisation}:
robust/probabilistic planning incorporates motion and uncertainty directly rather than
relying on PTV margins \cite{unkelbach2018}, which is the natural entry point for a
\emph{biomarker} error bar but has so far ingested mainly geometric and motion
uncertainty. Second, \textbf{technical validation gates clinical use}: imaging biomarkers
on an MR-linac require repeatability/reproducibility validation before driving decisions,
because the integrated scanner differs from a diagnostic system
\cite{vanhoudt2021qib}.% QUOTE-4 RESOLVED (vanhoudt2021qib) — verbatim re-pull, van Houdt et al. 2021, Eur J Cancer 153:64-71, Abstract (PMID 34144436; open access PMC8340311): "As the integrated MRI scanner differs from traditional MRI scanners, technical validation is an important aspect of this roadmap. We propose to integrate technical validation with clinical trials by the addition of a quality assurance procedure at the start of a trial, the acquisition of in vivo test-retest data to assess the repeatability, as well as a comparison between QIBs from MRIgRT systems and diagnostic MRI systems to assess the reproducibility." Paraphrase above is faithful; no number used. See CITATIONS.md.
{} That validation programme is now underway: quantitative $T_1$/$T_2$/ADC mapping is
feasible and accurate on the 1.5\,T MR-linac \cite{kooreman2019}, the Elekta Unity
consortium has issued ADC-acquisition recommendations for the integrated system
\cite{kooreman2020}, and a multi-institutional registry pools the clinical and technical
data needed to evaluate it \cite{demolvanotterloo2020}. \textbf{Takeaway:} the field is
converging on biomarker-driven adaptive dose, and the
two cautions define the demand it places on the quantitative-MRI side --- a calibrated,
technically validated error bar --- yet the propagation of \emph{biomarker uncertainty}
into the dose decision is the immature link. The rest of the review asks whether that
side can meet the demand: what is measured (Sec.~\ref{sec:foundations}), whether its
error bar is honest (Sec.~\ref{sec:trust}), and whether a trusted bar would change the
decision (Sec.~\ref{sec:voi}).

\section{Foundations: estimands and estimators}
\label{sec:foundations}

Before honesty can be assessed, the estimand and estimator must be fixed, because the
error bar a method reports is a property of the estimator, not of the biomarker alone.
The IVIM model was introduced to separate molecular diffusion from capillary-network
microcirculation within a voxel \cite{lebihan1986,lebihan1988}, treating randomly
oriented capillary flow as a pseudo-diffusion process \cite{lebihan2019}; dynamic
contrast-enhanced (DCE) MRI quantification rests on the compartmental transfer-constant
model \cite{tofts1991} and its standardised quantities $\Ktrans$, $v_e$, $k_{ep}$
\cite{tofts1999}.

The estimator is not innocent. Cram\'er--Rao/Fisher analysis shows the achievable
precision of IVIM $D$ and $f$ is strongly SNR-dependent \cite{zhang2013}, and the same
bounds can be inverted to design SNR-efficient acquisitions \cite{jalnefjord2019}.
Empirically, IVIM parameter estimates differ \emph{significantly} across fitting
algorithms in the upper abdomen, with Bayesian fitting giving the lowest inter-reader
and inter-subject variability \cite{barbieri2016}; simulation \cite{while2017} and
pancreatic-cancer data \cite{gurneychampion2018} corroborate the Bayesian advantage for
$D$ and $f$. A parallel line replaces hand-built fits with learned ones ---
neural-network IVIM/kurtosis mapping \cite{bertleff2017}, supervised and unsupervised
physics-informed deep fitting \cite{barbieri2020,kaandorp2021}, and self-supervised
forward-model estimation \cite{vasylechko2021} --- with recent work decomposing
predictive uncertainty into aleatoric and epistemic components \cite{casali2025} or
returning full posteriors that expose model degeneracy \cite{jallais2024}. DCE
shares the trajectory, e.g.\ Bayesian spatial tracer-kinetic estimation
\cite{schmid2006}. \textbf{Takeaway:} the uncertainty a method reports is inseparable
from the estimator that produced it, so the error bar an adaptive-RT decision consumes
must be judged per-estimator, not per-biomarker.

\section{Trust: is the error bar honest?}
\label{sec:trust}

\subsection{The general machinery of trust}
\label{sec:trust-machinery}

The UQ community supplies the toolkit the imaging literature draws on. Distribution-free
\textbf{conformal prediction} wraps any black-box predictor in a set/interval valid at a
user-chosen error rate \cite{shafer2008,angelopoulos2021}. \textbf{Calibration} methods
recalibrate regression intervals post-hoc \cite{kuleshov2018} and correct the systematic
over-confidence of modern deep networks \cite{guo2017}. Practical deep-UQ families ---
deep ensembles \cite{lakshminarayanan2017} and MC-dropout as approximate Bayesian
inference \cite{gal2016} --- give usable
uncertainty, but their quality \textbf{degrades under dataset shift}
\cite{ovadia2019} --- the operative warning for an MR-linac, whose integrated scanner is
itself a shift from the diagnostic systems on which such methods are trained.

\subsection{The metrology of trust in quantitative imaging}
\label{sec:trust-metrology}

Imaging has its own standards for honest reliability. The QIBA framework anchors
technical performance on three metrology areas --- linearity/bias, repeatability,
reproducibility \cite{raunig2015} --- with companion guidance on comparing algorithms
\cite{obuchowski2015} and on consistent terminology \cite{sullivan2015}. Reliability
reporting leans on the intraclass correlation coefficient
\cite{koo2016}% QUOTE-1 RESOLVED (koo2016) — verbatim re-pull, Koo & Li 2016, J Chiropr Med 15(2):155-163, Abstract/Conclusion (PMID 27330520): "This article provides a practical guideline for clinical researchers to choose the correct form of ICC and suggests the best practice of reporting ICC parameters in scientific publications." Paraphrase above is faithful; no number used. See CITATIONS.md.
{} and on limits-of-agreement rather than correlation for method comparison
\cite{blandaltman1986}.% QUOTE-2 RESOLVED (blandaltman1986) — verbatim re-pull, Bland & Altman 1986, Lancet 1(8476):307-310, Abstract (PMID 2868172): "Such investigations are often analysed inappropriately, notably by using correlation coefficients. The use of correlation is misleading. An alternative approach, based on graphical techniques and simple calculations, is described, together with the relation between this analysis and the assessment of repeatability." Paraphrase above is faithful; no number used. See CITATIONS.md.

\subsection{Where trust is empirically tested in body MRI}
\label{sec:trust-bodymri}

When these standards are applied to body diffusion/perfusion MRI, the error bars are
often \textbf{not} honest --- and the dishonesty is parameter- and cohort-dependent. ADC
is a comparatively reliable biomarker but its in-vivo reproducibility is markedly worse
than phantom reproducibility \cite{miquel2016}, and the body-DWI consensus already urged
baseline reproducibility studies as part of study design \cite{koh2007,padhani2009}. The
perfusion parameters are the weak point: short-term test--retest repeatability
coefficients for IVIM $f$ and $\Dstar$ reach ${\sim}126\%$ and ${\sim}197\%$ in rectal
cancer \cite{sun2017}, $\Dstar$ reproducibility is only moderate
($\mathrm{ICC}\approx0.55$, $\mathrm{CV}\approx20\%$) \cite{yang2019}, and multi-scanner
liver IVIM within-subject CVs run from ${\sim}5\%$ for $D$ up to ${\sim}14\%$ for
perfusion-fraction components \cite{simchick2024}. On the MR-linac itself, even
ADC --- the reliable end of the spectrum --- carries a sizeable test--retest repeatability
coefficient: ${\sim}17\%$ in rectal cancer \cite{eijkelenkamp2025} and
${\sim}13$--$31\%$ across head-and-neck targets \cite{habrich2022}, setting an explicit
floor below which an observed change is measurement noise, not response. Cross-modal
anchoring does not rescue $\Dstar$: its correlation with DCE $\Ktrans$ is weak in one rectal cohort ($r=0.389$)
\cite{sun2019} and non-significant in another \cite{yang2019}. \textbf{Takeaway:} trust
in body-MRI UQ is conditional --- broadly defensible for $D$/ADC, fragile for
$f$/$\Dstar$ --- yet most deep-learning UQ reports aggregate or marginal calibration
\cite{casali2025}, leaving the conditional failure region (exactly the regime an adaptive
dose change would rely on) under-characterised.

\section{Value of information: does a trusted bar change the dose decision?}
\label{sec:voi}

A calibrated interval earns its cost only if it moves an action --- here, a dose
decision. The decision-analytic literature supplies the apparatus: \textbf{decision curve
analysis / net benefit} evaluates a model by clinical net benefit across the decision
threshold \cite{vickers2006,vickers2016,vickers2019}, where the threshold encodes the
clinician's relative weighting of a missed disease against an unnecessary intervention;
performance frameworks argue such decision-analytic measures should be reported whenever a
model supports decisions \cite{steyerberg2010}. From the inference side,
\textbf{decision/loss-calibrated} methods tie the posterior approximation to the
downstream decision --- decision calibration \cite{zhao2021} and loss-calibrated
Bayesian inference \cite{vadera2021}. From health economics,
\textbf{value-of-information} analysis (EVPI) prices a measurement by combining the
probability it changes a decision with the benefit of that change \cite{felli1998},
under the decision-theoretic stance that choices follow expected net benefit rather than
statistical significance \cite{claxton1999}. \textbf{Takeaway:} the machinery to ask
``does this body-MRI interval change an oncologic dose decision?'' exists and is mature
--- but is rarely connected to quantitative-MRI uncertainty, which is precisely what an
adaptive-RT plan change would need.

\section{Gap map: open problems for biomarker-driven adaptive RT}
\label{sec:gaps}

The field's unresolved UQ-trust questions cluster at the \textbf{seams between the
axes} --- and, read for adaptive RT specifically, each is a problem that must be solved
before a biomarker error bar can be trusted to steer dose.

\begin{itemize}
  \item \textbf{G1 --- conditional honesty (trust seam).} Repeatability evidence shows
    $f$/$\Dstar$ are far less reliable than $D$/ADC
    \cite{sun2017,yang2019,simchick2024}, and even ADC carries a non-trivial on-device
    repeatability floor \cite{eijkelenkamp2025,habrich2022}, yet deep-learning IVIM-UQ
    typically reports marginal/aggregate calibration \cite{casali2025}. \emph{Open:} calibration
    conditioned on the regime where it fails, and conformal/coverage guarantees that
    respect it \cite{angelopoulos2021,ovadia2019} --- the regime an adaptive dose change
    would actually rely on.
  \item \textbf{G2 --- from calibration to dose decision (trust$\rightarrow$value seam).}
    The decision-analytic toolkit \cite{vickers2006,zhao2021} is rarely applied to
    quantitative-MRI uncertainty. \emph{Open:} when does a calibrated body-MRI interval
    actually change an oncologic dose decision versus merely being reported?
  \item \textbf{G3 --- uncertainty into dose (value$\rightarrow$action seam).}
    Adaptive-RT reviews call for technically validated biomarkers to steer dose
    \cite{otazo2021,vanhoudt2021qib}, and robust planning can ingest uncertainty
    \cite{unkelbach2018}. \emph{Open:} end-to-end propagation of \emph{biomarker}
    uncertainty (not just geometric/motion uncertainty) into the adaptive dose decision,
    and the identifiability walls that cap it.
  \item \textbf{G4 --- estimator-dependence of trust (cross-cutting).} Reported
    uncertainty depends on the estimator \cite{barbieri2016,barbieri2020}. \emph{Open:}
    standards that make UQ comparable across classical and learned estimators, extending
    QIBA metrology \cite{raunig2015,sullivan2015} to the uncertainty itself, not only the
    point estimate --- a prerequisite for trusting any of them on an MR-linac.
\end{itemize}

\section{Discussion and limitations}
\label{sec:scope-honest}

This review establishes three things for the biomarker-to-dose problem. First, trust in
body-MRI uncertainty is \emph{conditional}: the diffusion coefficient $D$ and ADC carry
broadly defensible error bars, whereas the perfusion parameters $f$ and $\Dstar$ do not,
and most deep-learning UQ reports only marginal, aggregate calibration that does not
characterise where the failure concentrates (Sec.~\ref{sec:trust}). Second, the
decision-analytic apparatus needed to price a calibrated interval against a dose
decision --- net benefit, decision calibration, value of information --- is mature, but
is rarely connected to quantitative-MRI uncertainty (Sec.~\ref{sec:voi}). Third, the
propagation of \emph{biomarker} uncertainty into the adaptive dose decision, as opposed
to geometric or motion uncertainty, remains the immature link (Sec.~\ref{sec:action}).
The four open problems (Sec.~\ref{sec:gaps}) sit at exactly these seams.

The review deliberately draws a scope boundary. It covers uncertainty quantification for
quantitative and diffusion body MRI and its decision-use in MR-guided adaptive
radiotherapy; it excludes non-quantitative MRI, neuro-only work, generic computer-vision
UQ with no quantitative-MRI tie, and radiotherapy that is not MR-guided. That boundary
keeps the survey on the dose-decision link and is recorded in the repository's
\texttt{ASSUMPTIONS.md}.

The principal limitation is methodological: this is a single-author \emph{narrative}
synthesis, not a systematic review. It follows no pre-registered search protocol, and the
selection of works --- though chosen to span the estimation, calibration,
decision-analytic, and MR-guided-RT literatures the problem touches --- is expert-curated
and therefore subject to selection bias. A systematic review with an explicit
search-and-inclusion protocol, ideally with more than one reviewer, would be the natural
next step to harden the coverage claims made here. Every empirical statement is
nonetheless attributed to a verified external source.

% --- phiRO / Elsevier declaration sections (unnumbered) ---
\section*{Declaration of competing interest}
The author declares no competing interests.

\section*{Funding}
This research received no specific grant from any funding agency in the public,
commercial, or not-for-profit sectors; it was conducted independently by the author.

\section*{CRediT authorship contribution statement}
\textbf{Avery Karlin:} Conceptualization, Methodology, Software, Formal analysis,
Investigation, Data curation, Writing -- original draft, Writing -- review and editing.

\section*{Data availability}
This review reports no new data. The verified citation base and the script that checks
that every reference resolves are openly available in the project repository.

\section*{Declaration of generative AI and AI-assisted technologies in the manuscript preparation process}
During the preparation of this work the author used Claude (Anthropic), including Claude
Code, for code scaffolding, manuscript drafting assistance, and copy-editing. After using
these tools the author reviewed and edited the content as needed and takes full
responsibility for the content of the publication.

\bibliographystyle{elsarticle-num}
\bibliography{limbo}

\end{document}
